\section{Spectral and convexity preliminaries}\label{pre:section}

Two elementary facts let us work with measurable potentials and identify
the equality condition in the potential estimate. The first is the form
argument used in the matrix inequality~\cite[Lemma~7.1]{MLT}; we give its
proof to fix the precise regularity available below.

\begin{lemma}\label{pre:spectral}
Let $p>1$ and let $W:\R\to\C^{m\times m}$ be measurable, Hermitian and
positive semidefinite, with $\int_\R\tr(W^p)<\infty$.
The form $h_W$ on $H^1(\R;\C^m)$ is closed and bounded below.
Multiplication by $W^{1/2}$ is compact from $H^1$ to $L^2$.
The negative spectrum of $H_W$ consists of isolated eigenvalues of finite
multiplicity, with possible accumulation only at zero. Its eigenfunctions
belong to $H^1$ and satisfy $-\psi''-W\psi=\lambda\psi$ distributionally.
\end{lemma}

\begin{proof}
Set $w=\|W\|_{\op}$. Then $w^p\le\tr(W^p)$, so $w\in L^p$.
The one-dimensional Sobolev estimate, applied to the norm of a vector
function, gives $\|\psi\|_\infty^2\le2\|\psi\|_2\|\psi'\|_2$.
H\"older's inequality, interpolation and Young's inequality imply
\begin{equation}\label{pre:form-bound}
 \int_\R\langle\psi,W\psi\rangle
 \le 2^{1/p}\|w\|_p\|\psi'\|_2^{1/p}\|\psi\|_2^{2-1/p}
 \le \varepsilon\|\psi'\|_2^2+C_\varepsilon\|\psi\|_2^2.
\end{equation}
A sufficiently shifted form norm is therefore equivalent to the $H^1$
norm. This proves closedness and semiboundedness, and polarization gives
continuity of the potential form on $H^1$.

Write $T\psi=W^{1/2}\psi$. The truncated multiplier
$T_{R,L}=\mathbf1_{\{|x|\le R,\,w\le L\}}T$ is compact by the compact
embedding of $H^1([-R,R])$ into $L^2([-R,R])$. Moreover,
\[
 \|(T-T_{R,L})\psi\|_2^2
 \le C\|w\mathbf1_{\{|x|>R\}\cup\{w>L\}}\|_p\|\psi\|_{H^1}^2.
\]
The coefficient tends to zero, so $T$ is compact.
For each $a>0$, an infinite-dimensional spectral subspace in
$(-\infty,-a]$ would contain an $L^2$-orthonormal sequence bounded in
$H^1$ by~\eqref{pre:form-bound}. A subsequence converges weakly to zero
in $H^1$, hence its images under $T$ converge strongly to zero. This
contradicts
$-a\ge h_W[\psi]\ge-\|T\psi\|_2^2$ for unit vectors in that subspace.
Thus every such subspace is finite-dimensional. Finally, the weak form
equation tested against smooth compactly supported functions is the stated
distributional equation; its potential term is locally integrable because
$W\in L^p_{\mathrm{loc}}$ and $\psi$ is locally bounded.
\end{proof}

\begin{lemma}[Trace Young inequality and its equality condition]\label{pre:young}
For positive semidefinite matrices $W,Z$ of the same finite size, set
\begin{equation}\label{pre:young-defect}
 Y(W,Z)=D_\sigma\tr(W^p)-\tr(WZ)+\tr(Z^q).
\end{equation}
Then $Y(W,Z)\ge0$, with equality if and only if $W=qZ^r$.
Also,
\begin{equation}\label{pre:young-coercive}
 Y(W,Z)\ge\tfrac12\tr(Z^q)
       -(2^{p-1}-1)D_\sigma\tr(W^p).
\end{equation}
\end{lemma}

\begin{proof}
Scalar maximization over $z\ge0$ gives
$yz-z^q\le D_\sigma y^p$ for $y\ge0$, with equality precisely when
$y=qz^r$. In an orthonormal eigenbasis $(v_j)$ of $Z$, let $z_j$ be its
eigenvalues and put $y_j=\langle v_j,Wv_j\rangle$. Strict convexity of
$t^p$ on $[0,\infty)$ gives
\[
 \tr(WZ)-\tr(Z^q)
 \le D_\sigma\sum_j y_j^p
 \le D_\sigma\sum_j\langle v_j,W^pv_j\rangle.
\]
Equality in the second inequality holds exactly when each $v_j$ belongs
to an eigenspace of $W$. Equality in the first then requires
$Wv_j=qz_j^rv_j$ for every $j$. This proves both directions of the
equality assertion, including zero eigenvalues. Applying the same
inequality to $2W,Z$ gives
$\tr(WZ)\le\tfrac12\tr(Z^q)+2^{p-1}D_\sigma\tr(W^p)$,
which is~\eqref{pre:young-coercive}.
\end{proof}
